\cleardoublepage
\chapternonum{摘要}

端粒到端粒（telomere-to-telomere, T2T）基因组的出现，使植物复杂重复区、着丝粒和端粒等长期难以解析的区域得以完整表征，为系统研究植物基因组结构与演化提供了新契机。然而，当前植物T2T基因组仍存在资源分散、跨物种比较不足、复杂区域研究有限及组装技术路线缺乏系统评估等问题。围绕上述问题，本论文以植物T2T基因组为对象，系统开展了资源整合、比较基因组学分析、复杂区域演化解析等方面的研究。

本论文首先收集了230个植物T2T基因组，对缺失编码基因注释的基因组进行了从头注释、质量评估。并开发了iSeq和iOmics分析工具、PlanT2T数据库和T2T-Hub在线分析平台。在此基础上，从整体序列层面分析了GC含量演化、启动子与转录因子的协同进化以及NLR基因家族的演化规律。并基于116个物种构建时间校准的系统发育树，定位了多个关键全基因组加倍（WGD）事件，并揭示了3亿年来WGD与基因家族扩张与环境变化之间的关系，并将相关结论扩展到1058个非T2T基因组中进行了验证。本论文还开展了重复序列的\textit{de novo}注释，阐明了LTR反转录转座子在植物基因组扩张中的核心作用，并揭示了重复序列与调控元件互作及性状变异之间的联系。然后，本论文对植物着丝粒的结构组成、表观特征和演化模式进行了系统比较，提出了TE-TR循环演化模型。在方法学方面，本论文对植物T2T基因组的调查、初步组装、挂载和补洞策略进行了多平台、跨物种、不同倍性的系统基准测试。结果表明Hifiasm、HapHiC和TGS-GapCloser分别在组装、挂载和补洞阶段表现出最稳定的综合优势。基于该结论，本论文完成了同源六倍体落地生根的T2T基因组组装。最后，通过比较T2T与非T2T参考基因组在ATAC-seq、RNA-seq和BS-seq中的表现，证明了T2T基因组在多组学信号识别和复杂区域信息恢复方面具有显著优势。

总之，本论文系统揭示了植物T2T基因组的结构特征及其演化规律，建立了适用于复杂植物基因组的T2T组装技术路线，并验证了T2T参考基因组在多组学研究中的应用价值，为植物基因组学研究和相关资源平台建设提供了新的数据基础与方法基础。

\noindent\textbf{关键词：}植物；T2T基因组；演化；重复序列；着丝粒

\cleardoublepage
\chapternonum{Abstract}

The emergence of telomere-to-telomere (T2T) genomes has enabled the complete characterization of long-intractable genomic regions in plants, including complex repeat arrays, centromeres, and telomeres, offering new opportunities to systematically investigate plant genome structure and evolution. However, current plant T2T genome research still faces several challenges, including dispersed data resources, insufficient cross-species comparisons, limited exploration of complex genomic regions, and a lack of systematic evaluation of assembly methodologies. To address these issues, this dissertation focuses on plant T2T genomes and conducts systematic studies in resource integration, comparative genomics, and evolutionary analysis of complex regions.

First, this dissertation collected 230 plant T2T genomes, performed de novo annotation and quality assessment for those lacking protein-coding gene annotations, and developed two analytical tools (iSeq and iOmics), a comprehensive database (PlanT2T), and an online analysis platform (T2T-Hub). Based on these resources, we analyzed the evolution of GC content, the co-evolution of promoters and transcription factors, and the evolutionary patterns of the NLR gene family at the whole-genome level. Using a time-calibrated phylogenetic tree constructed from 116 species, we identified multiple key whole-genome duplication (WGD) events and revealed the relationships between WGD, gene family expansion, and environmental changes over the past 300 million years. These findings were further validated in 1,058 non-T2T genomes. We also performed de novo repeat annotation, elucidated the central role of LTR retrotransposons in plant genome expansion, and uncovered links between repeat elements, regulatory element interactions, and trait variation. Next, this dissertation systematically compared the structural composition, epigenetic features, and evolutionary patterns of plant centromeres, proposing a TE-TR (transposable element–tandem repeat) cyclic evolution model. In terms of methodology, we conducted a multi-platform, cross-species, and ploidy-aware systematic benchmark of strategies for plant T2T genome survey, initial assembly, scaffolding, and gap filling. Our results demonstrated that Hifiasm, HapHiC, and TGS-GapCloser exhibited the most consistently superior performance in assembly, scaffolding, and gap filling, respectively. Based on these findings, we assembled a T2T genome for the allohexaploid \textit{Kalanchoe laetivirens}. Finally, by comparing T2T and non-T2T reference genomes in ATAC-seq, RNA-seq, and BS-seq analyses, we demonstrated the significant advantages of T2T genomes in multi-omics signal detection and information recovery from complex genomic regions.

In summary, this dissertation systematically reveals the structural characteristics and evolutionary patterns of plant T2T genomes, establishes a T2T assembly pipeline suitable for complex plant genomes, and validates the utility of T2T reference genomes in multi-omics studies. These findings provide new data resources and methodological foundations for plant genomics research and the development of related bioinformatics platforms.

\noindent\textbf{Keywords:} plant; T2T genome; evolution; repetitive sequences; centromere
